Any claim about change over time needs dates, and \DatedById{} of the
\NBooks{} book identifiers carry one, because they were filed as
\texttt{AUTHOR\_YEAR\_Title}. The rest --- typed editions and rescans filed
under a plain title --- carry none, and dropping them would have discarded a
third of the corpus.

Rather than supply dates from outside the collection, we read them out of it.
Title pages print the year, in arabic or roman numerals, and they sit on the
first few leaves. A first attempt simply took the earliest plausible year in
the front matter, and it failed in a way worth recording, because it is the
failure mode a digitised library invites. It dated Shaftesbury's
\emph{Characteristicks} to 1671 --- his birth year, printed in a modern
editor's note --- and gave four different works by Hutcheson the same year,
1740, because all four are reproductions of a modern edition whose preface
carries that date. Modern apparatus sits in front of the original title page
and prints dates of its own, and the naive reader cannot tell the difference
\citep{cordell2017}.

The version used here accepts a year only when the page it sits on looks like
a period title page, which in practice means an imprint line naming a place or
a printer, or when the year is set in roman numerals, which modern front
matter does not use. Any candidate equal to the author's own birth or death
year, as recorded in the project's author registry, is refused outright. That
refused \DatesRefused{} candidates and accepted \DatesRecovered{}, raising
coverage from \CoverageBefore\% to \CoverageAfter\% of passages;
\StillUndated{} works remain undated and are excluded from every analysis over
time rather than guessed at.

The method can be checked, because the books whose identifiers already carry a
year also have title pages. Running the same reader over them gives
\DateCheckN{} readable cases, of which \DateCheckExact{} match the catalogue
year exactly (\DateCheckExactPct\%) and \DateCheckWithinOne{} fall within a
year. That is a small validation set and we do not claim more from it than it
can carry: the reader is precise when it fires and fires rarely, which is the
behaviour we wanted.

One caveat governs everything downstream. What a title page prints is the year
of \emph{that impression}, not of first publication. A 1755 reprint of a 1725
argument is dated here to 1755. For the vocabulary curves in
\S\ref{sec:diachrony} this is arguably the right quantity --- it dates the
copy a reader could have held --- but it compresses the early period and must
not be read as a chronology of composition.
